% Full provenance: CONTRAANGULO_ELECTRON_BARBERO_CATALAN.md, §§2--10.
% Editorial records, history and source locators are preserved in gestion/.
\section{Composition and reconstruction of the angular, electronic and constitutive readouts}
\label{md:compos:seccion}

\subsection{Action coefficient and construction of the counterangle}
\label{md:compos:contraangulo}

The functional abbreviation $B(\alpha)$ collects an evaluation at $\alpha$;
its argument belongs to the function and does not denote the product of an
independent constant and $\alpha$. To keep the operations involved in this
abbreviation explicit, write
\begin{equation}
 B(\alpha)=\frac{9}{16}\alpha^2-\frac{5}{9}\alpha^3
 +\frac{7}{48}\alpha^4-\frac{1}{54}\alpha^5
 -\frac{90}{\pi}\mathcal C_\pi(\alpha),
 \label{md:compos:b-alpha}
\end{equation}
\begin{equation}
 \mathcal C_\pi(\alpha)=169\alpha^6+
 \frac{D_A\alpha^7}{1-D_A\alpha},
 \qquad D_A=\exp\!\left(-\frac{100\pi\alpha}{9}\right).
 \label{md:compos:correccion-regional}
\end{equation}
The complete identity is
\begin{equation}
 C^*_{\deg}=\sqrt{\frac{1000\alpha}{\varphi}}+2B(\alpha).
 \label{md:compos:contraangulo-abreviado}
\end{equation}
In particular, $\alpha$ belongs to the radicand, and the radical term is
retained together with $2B(\alpha)$. This identity rewrites the map
constructed earlier while preserving its object and genealogy.

The original definition keeps the dimensionless coefficient separate from
the action section:
\begin{equation}
 \eta_{\mathrm{ret}}=H_5-\frac{90}{\pi}\mathcal C_\pi,
 \qquad
 \hbar_{\mathrm{ret}}=\eta_{\mathrm{ret}}\,10^{-34}\mathcal U_S.
 \label{md:compos:seccion-accion}
\end{equation}
Hence,
\begin{equation}
 \boxed{C^*_{\deg}=2\left(
 \frac{\hbar_{\mathrm{ret}}}{10^{-34}\mathcal U_S}+\alpha\right)
 =2(\eta_{\mathrm{ret}}+\alpha).}
 \label{md:compos:contraangulo-accion}
\end{equation}
The expression uses the dimensionless coefficient of the reduced Planck
section. The constant $h$ is subsequently obtained through $h=2\pi\hbar$.
The composition extracts the action coefficient, combines it with $\alpha$
and applies bidirectional closure. This preserves the dimensional type of
each operation.

The mass construction expresses particle and antiparticle as two oriented
sections of the same genealogical band. The dodecaphasic operation is
\begin{equation}
 C_M(\tau)=-\operatorname{rev}(\tau),
 \qquad \tau_e=+++---+++---,
 \qquad C_M(\tau_e)=\tau_e,
 \label{md:compos:conjugacion-dodecafasica}
\end{equation}
extended to the register and the charge representation. The central route
is fixed under $C_M$, whereas the charge representation distinguishes the
electron from the positron. This composition gives precise content to the
bidirectional readout. The angular factor of two retains its closure
function; physical conjugation additionally retains the stated operators.
The return at $216$ restores the oriented lift: its structure does not
consist in assigning $108$ updates to the particle and another $108$ to
the antiparticle.

\subsubsection{Radial reconstruction of the action coefficient}
\label{md:compos:quinientos}

Write $q_{\mathrm{rad}}=q_\varepsilon$ and $\epsilon=\varepsilon$ for the
radial ratio and orientation in \eqref{rec:eq:eta-orientada}.
The oriented radial invariant $\chi_{\mathrm{rad}}=\epsilon\log q_{\mathrm{rad}}$
allows the previously constructed reconstruction to be written as
\begin{equation}
 \frac{\eta_{\mathrm{ret}}}{\alpha}
 =-500\tanh\!\left(\frac{\chi_{\mathrm{rad}}}{2}\right)-1.
 \label{md:compos:restitucion-radial}
\end{equation}
With $\xi=C^*_{\deg}/A_{\deg}$ and $A_{\deg}=1000\alpha$, one also obtains
\begin{equation}
 \eta_{\mathrm{ret}}=\alpha(500\xi-1).
 \label{md:compos:restitucion-xi}
\end{equation}
These are downstream reconstruction maps for the same section; its
construction remains in $H_5$ and in the regional return.

\subsection{Angular charts and conservation of action}
\label{md:compos:cartas}

The degree chart has a full turn of $360$ units and preserves the
partitions $12\cdot30=4\cdot90=3\cdot120=360$. The constructed angular pair is
\begin{equation}
 A_{\deg}=1000\alpha,
 \qquad C^*_{\deg}=2(\eta_{\mathrm{ret}}+\alpha).
 \label{md:compos:par-angular}
\end{equation}
Conversion is performed once:
\begin{equation}
 x=\frac{\pi}{180}A_{\deg},\qquad
 y=\frac{\pi}{180}C^*_{\deg},\qquad
 q_\pm=\exp(-x\mp y).
 \label{md:compos:canales}
\end{equation}
The exponents use lifted real representatives and retain the number of
turns. The action satisfies
\begin{equation}
 \hbar\theta_{\mathrm{rad}}
 =\frac{h\theta_{\deg}}{360}.
 \label{md:compos:accion-cartas}
\end{equation}
For the bilinear Barbero functional, chart transport gives
\begin{equation}
 \frac{\pi A_{\deg}C^*_{\deg}}{180}
 =\frac{180xy}{\pi}.
 \label{md:compos:bilineal-cartas}
\end{equation}
The bilinear transformation is part of the same conversion applied to the
exponential arguments; omitting it would change the functional.

\subsection{Complete electronic readout and action return}
\label{md:compos:electron}

The two stages of the same electronic coordinate are retained:
\begin{equation}
 \varepsilon_e^{(0)}=\frac{\sqrt{3}}{4}
 \left[\exp\!\left(\frac{\varphi}{\pi^2}\right)-22\alpha^3\right]
 (1+15\Delta_4),\qquad
 \Delta_4=\frac{\pi-e\log\pi}{270}
 \left(1+\frac{\pi}{729}\right),
 \label{md:compos:electron-basal}
\end{equation}
\begin{equation}
 \boxed{\varepsilon_e^\beta=\varepsilon_e^{(0)}
 \left(\frac{H_5}{\eta_{\mathrm{ret}}}\right)^{80-54\alpha+6\Delta_4}.}
 \label{md:compos:electron-completo}
\end{equation}
The centre $(5,5)$, its inversion, the deformations $(8,2)/(2,8)$ and the
unit quotient defect precede the value. The norm $\sqrt{3}/4$ comes from the
commutator between additive and multiplicative projections; the exponential
character uses the orientation $\varphi/\pi^2$; $22=27-5$ retains a regional
complement and $15=3\cdot5$ the corresponding incidences. The route readout
maps produce $(80,54,6)$. The ratio $H_5/\eta_{\mathrm{ret}}$ incorporates
the action return into the same central section.

The role of the counterangle is made explicit by substituting its proven
inverse:
\begin{equation}
 \boxed{\varepsilon_e^\beta=\varepsilon_e^{(0)}
 \left(\frac{H_5}{C^*_{\deg}/2-\alpha}\right)^{80-54\alpha+6\Delta_4}.}
 \label{md:compos:electron-contraangulo}
\end{equation}
This expression preserves the preceding electronic generation. The same
action section that participates in the counterangle determines the factor
taking the basal readout to the complete one. The ratio cancels the common
dimensional basis; the absolute mass realization retains its temporal
readout operator and dimensional chart. The substitution composes constructed
dependencies without introducing an external angle as an electronic generator.

\subsection{Barbero functional and the hexad--octad transition}
\label{md:compos:barbero}

The incidence preserves the same marked pair as the support is enlarged:
\begin{equation}
 \mathcal F_6=\mathsf H\times\binom{\mathsf H}{2}
 \hookrightarrow
 \mathcal F_8=\mathsf O\times\binom{\mathsf H}{2},
 \qquad |\mathcal F_6|=90,\qquad |\mathcal F_8|=120.
 \label{md:compos:incidencia}
\end{equation}
The cardinalities are $6\cdot15$ and $8\cdot15$; the shared reference allows
the two incidences to be compared. The normalized difference is $-1/12$.
The dodecaphasic chain retains twelve hexadic contributions and one octadic
boundary contribution with the opposite sign. The readout operator
\begin{equation}
 S_m(q)=\frac{q^m}{1-q^{3m}}
 =\sum_{j\geq0}q^{m+3mj}
 \label{md:compos:serie-retorno}
\end{equation}
retains height and returns. This yields $12S_{90}-S_{120}$ and its
difference between the two orientations. The complete functional is
\begin{equation}
 \gamma=\frac{\pi A_{\deg}C^*_{\deg}}{180}
 +12\bigl[S_{90}(q_-)-S_{90}(q_+)\bigr]
 -\bigl[S_{120}(q_-)-S_{120}(q_+)\bigr].
 \label{md:compos:funcional-barbero}
\end{equation}
In this realization, Barbero is the joint oriented evaluation of the
angular pair and the returns of the incidence transition. Its participation
in the area scale is subsequently composed with $\ell_P^2=G\hbar/c^3$ and
with the occupation operator. Entanglement entropy additionally retains the
reduced state and its probabilities; its reconstruction requires those data
and is not determined by scalar area alone.

\subsubsection{Elimination of the common electron--Barbero coordinate}
\label{md:compos:eliminacion}

Define $\kappa_e=80-54\alpha+6\Delta_4$. On the positive branch, with
$\kappa_e\neq0$, the electronic equation implies exactly
\begin{equation}
 \eta_{\mathrm{ret}}=H_5
 \left(\frac{\varepsilon_e^{(0)}}{\varepsilon_e^\beta}\right)^{1/\kappa_e}.
 \label{md:compos:eliminacion-eta}
\end{equation}
To express the composition while keeping its operations explicit, put
$x=50\pi\alpha/9$, $y=\pi(\eta_{\mathrm{ret}}+\alpha)/90$ and
$\mathcal R(q)=12S_{90}(q)-S_{120}(q)$. Then
\begin{equation}
 \gamma=\frac{100\pi}{9}\alpha(\eta_{\mathrm{ret}}+\alpha)
 +\mathcal R(e^{-x+y})-\mathcal R(e^{-x-y}).
 \label{md:compos:barbero-eta}
\end{equation}
Substitution of \eqref{md:compos:eliminacion-eta} into this identity yields
an explicit relation between the complete electronic readout and Barbero,
retaining $\alpha$, $H_5$, $\Delta_4$ and the incidences. It is a
compatibility consequence of the construction. The angular physical branch
used here satisfies $0<y<x$; the operation introduces no metrological mass
as generating data.

In the extension with orientation label $\epsilon=\pm1$,
$C^*_{\deg,\epsilon}=2\epsilon(\eta_{\mathrm{ret}}+\alpha)$, the exchange
$\epsilon\mapsto-\epsilon$ fixes $\eta_{\mathrm{ret}}$, preserves the
electronic expression and interchanges $q_+$ and $q_-$. The Barbero
functional changes sign. This covariance distinguishes a scalar action
readout from an oriented evaluation. Identification with charge conjugation
additionally retains $C_M$ and the electromagnetic sheet representation;
the angular sign alone does not replace these operators.

\subsection{Oriented Catalan moment and constitutive response}
\label{md:compos:catalan}

The dodecaphasic cycle realizes an oriented quarter-turn plane. Its
resolvent coefficient is $k_4(s)=s/(1+s^2)$. The depth-two moment is fixed by
\begin{equation}
 \mathcal C_2(s)=\sum_{k\geq0}
 \frac{(-1)^ks^{2k+1}}{(2k+1)^2},\qquad
 \mathcal D=s\frac{d}{ds},\qquad
 \mathcal D^2\mathcal C_2(s)=k_4(s),\qquad
 \mathcal C_2(1)=G_{\mathrm{Cat}}.
 \label{md:compos:momento-catalan}
\end{equation}
The differential identity is used for $0<s<1$; the endpoint value is
understood with the previously established convergence and moment
reconstruction conditions. The intrinsic definition of Catalan is the
endpoint value of this oriented depth moment. Its collection retains more
information than the endpoint value. The constitutive response uses that
collection:
\begin{equation}
 f(s)=\frac{1+s+s^2}{s(1+s)}\mathcal D^2\mathcal C_2(s)
 =\frac{1-s^3}{1-s^4},\qquad
 r_\pm=f(s_\pm),\qquad s_\pm=q_\pm^{30}.
 \label{md:compos:catalan-respuesta}
\end{equation}
The powers $3$ and $4$ correspond to $90=3\cdot30$ and $120=4\cdot30$.
\begin{proposition}[Functional reconstruction with a boundary condition]
\label{md:compos:inversa-catalan}
The constitutive collection $f$ determines the oriented moment through
\[
 \mathcal C_2(s)=\int_0^s\log\frac{s}{t}\,
 \frac{1+t}{1+t+t^2}f(t)\,dt,\qquad 0<s<1.
\]
It is the unique solution of
$\mathcal D^2\mathcal C_2(s)=s(1+s)f(s)/(1+s+s^2)$
that tends to zero at the origin.
\end{proposition}
\begin{proof}
For the constructed collection,
$(1+t)f(t)/(1+t+t^2)=1/(1+t^2)$. The integral converges, and its
logarithmic derivative is
$\mathcal D\mathcal C_2(s)=\int_0^s(1+t^2)^{-1}\,dt$.
Applying $\mathcal D$ again yields $s/(1+s^2)$. The integral is
$O(s)$ at the origin. Two solutions differ by $a+b\log s$;
the zero limit forces $a=b=0$.
For $s<1$, integrating the geometric series on compact subsets and using
$\int_0^s t^{2k}\log(s/t)\,dt=s^{2k+1}/(2k+1)^2$
recovers the series in \eqref{md:compos:momento-catalan}.
Its dominated absolute convergence allows $s\uparrow1$ and recovers
Catalan's constant.
\end{proof}
The boundary condition fixes the two integration constants.
The reconstruction uses the constitutive function together with its
domain; the endpoint value of Catalan's constant and the two channel
values are readouts of that already constructed function. At the endpoint,
the twice-differentiated identity retains its interpretation through an
Abel limit.

The constitutive operator $V=r_+P_++r_-P_-$ determines
\begin{equation}
 \widehat\varepsilon=r_+^2,\qquad \widehat\mu=r_-^2,
 \qquad \widehat Z=\frac{r_-}{r_+},\qquad
 \widehat c=(r_+r_-)^{-1}.
 \label{md:compos:lectores-constitutivos}
\end{equation}
With $\psi=f^{-1}$ and
\begin{equation}
 \mathcal H_B(s)=\frac{12s^3}{1-s^9}
 -\frac{s^4}{1-s^{12}}-\frac{(\log s)^2}{20\pi},
 \label{md:compos:h-barbero}
\end{equation}
the constructed factorization takes the form
\begin{equation}
 \gamma=\mathcal H_B(\psi(r_-))-\mathcal H_B(\psi(r_+)).
 \label{md:compos:factorizacion-barbero}
\end{equation}
The logarithmic part reconstructs the angular contribution to Barbero
exactly, retained in degrees or transported in full to radians.

\subsection{Reconstruction from normalized speed and Barbero}
\label{md:compos:difeomorfismo}

\paragraph{Proposition.}
In the normalized constitutive collection of two labelled channels, the map
\begin{equation}
 (r_+,r_-)\longmapsto(\widehat c,\gamma),\qquad
 (3/4,1)^2\longrightarrow(1,16/9)\times\mathbb R
 \label{md:compos:mapa-global}
\end{equation}
is a diffeomorphism. Its inverse recovers the two labelled eigenvalues.
The sheet projectors belong to the initial representation: two scalars do
not recover an arbitrary spatial basis.

\paragraph{Proof.}
The regular expression
$f(s)=(1+s+s^2)/(1+s+s^2+s^3)$ gives
\begin{equation}
 f'(s)=-\frac{s^2(s^2+2s+3)}{(1+s+s^2+s^3)^2}<0.
 \label{md:compos:derivada-f}
\end{equation}
Thus $f$ is analytic and strictly decreasing, with $f(0^+)=1$ and
$f(1^-)=3/4$. For $0<s<1$,
\begin{equation}
 \mathcal H_B'(s)=
 \frac{36s^2(1+2s^9)}{(1-s^9)^2}
 -\frac{4s^3(1+2s^{12})}{(1-s^{12})^2}
 -\frac{\log s}{10\pi s}>0.
 \label{md:compos:derivada-h}
\end{equation}
Indeed, the ratio of the first positive term to the second positive term
before its subtraction is
\begin{equation}
 \frac{9}{s}\frac{1+2s^9}{1+2s^{12}}
 \left(\frac{1-s^{12}}{1-s^9}\right)^2>9.
 \label{md:compos:cota-derivada-h}
\end{equation}
The last summand of \eqref{md:compos:derivada-h} is also positive.
Hence $F_B=\mathcal H_B\circ\psi$ has strictly negative derivative,
$F_B(3/4^+)=+\infty$ and $F_B(1^-)=-\infty$.

Fix $c=\widehat c$, write $z=\log\widehat Z$ and
$r_\pm=c^{-1/2}\exp(\mp z/2)$. The admissible interval is
\begin{equation}
 |z|<a(c):=\min\left\{\log c,\log\frac{16}{9c}\right\}.
 \label{md:compos:dominio-z}
\end{equation}
The function $\Gamma_c(z)=F_B(r_-)-F_B(r_+)$ satisfies
\begin{equation}
 \Gamma_c'(z)=\frac12
 \bigl[r_-F_B'(r_-)+r_+F_B'(r_+)\bigr]<0.
 \label{md:compos:derivada-gamma}
\end{equation}
As $z$ increases to $a(c)$, either $r_-$ reaches $1$, or $r_+$ reaches $3/4$,
or both occur; $\Gamma_c$ tends to $-\infty$. At the opposite endpoint it
tends to $+\infty$. There is therefore a unique solution for every real
$\gamma$. The nonzero derivative and the implicit function theorem provide
a smooth global inverse. This proves the statement. In the original
orientation $y>0$, one has $z<0$ and $\gamma>0$.

Normalized speed retains the product of the responses; Barbero
monotonically distinguishes their oriented distribution. Together they
reconstruct the spectral information that neither readout determines alone.

The statement concerns the collection of constitutive readout operators. The image
actually produced by the full HMT generator is a subset of that collection:
restriction to that image preserves injectivity and causal order.
Normalization and sheet markings are necessary data for reconstruction.
The proof neither introduces CODATA values to select channels nor asserts
that every synthetic point in the collection is a generated state.

\subsection{Reconstruction on the generated image}
\label{md:compos:restitucion-completa}

Once $r_\pm$ have been recovered from $(\widehat c,\gamma)$, one obtains
$s_\pm=\psi(r_\pm)$. Then
\begin{equation}
 A_{\deg}=-\frac{3}{\pi}\log(s_+s_-),\qquad
 C^*_{\deg}=\frac{3}{\pi}\log\frac{s_-}{s_+}.
 \label{md:compos:restitucion-angular}
\end{equation}
On the generated image and in its positive orientation,
\begin{equation}
 \alpha=\frac{A_{\deg}}{1000},\qquad
 \eta_{\mathrm{ret}}=\frac{C^*_{\deg}}{2}-\alpha>0.
 \label{md:compos:restitucion-alpha-eta}
\end{equation}
With the regional coordinates $\pi,\varphi,e$ and the dimensional sections
already constructed, this reconstruction recovers $H_5$, the factor
$R_{\mathrm{act}}=H_5/\eta_{\mathrm{ret}}$ and the complete electronic
readout of Section~\ref{md:compos:electron}. It also recovers the Planck
section. If the twelve-block positional base, origin and carry are retained
as well, the identity
\begin{equation}
 \kappa_{\mathrm{ag}}=\pi+e-\varphi-4-\alpha
 \label{md:compos:restitucion-k}
\end{equation}
allows reconstruction of the register $K$ described by its readout operator.
This is a reconstruction path; its reverse direction does not constitute
a new generating arrow from external physical values.

Gravitation and thermal conversion retain the sections of the same
state. Theorem~\ref{md:rigidez:teorema} and
\eqref{md:rigidez:reloj} first determine $L_\alpha$ and its
circular representation $R_{108}=L_\alpha$:
\begin{equation}
 G_a=\frac{c_{\mathrm{int}}^3R_{108}^2}{\hbar_a},\qquad
 \frac{G_a\hbar_a}{c_{\mathrm{int}}^3}=R_{108}^2,
 \qquad
 k_B\Theta_{\mathrm{clk},a}\log3=\frac{h_a}{108t_0}.
 \label{md:compos:gravedad-termica}
\end{equation}
The chronological radius, physical speed and thermal basis retain their
construction maps. In particular, the clock of the circular representation
is tied to $L_\alpha$ and $c_{\mathrm{int}}$. The first
composition makes action--gravitation reciprocity explicit; the second
combines action, duration and tritic information, with temperature and
$k_B$ individually specified in their declared bases.

\subsection{Structural definitions of the composite readouts}
\label{md:compos:definiciones}

\begin{description}
 \item[$\alpha$.] Scalar coordinate of dodecaphasic closure compatible with
 the regions and the arithmetic--geometric register $K$; it transmits that
 compatibility to the angular and action sectors.
 \item[Counterangle.] Angular coordinate of bidirectional closure combining
 the coefficient of the return action section with $\alpha$ in the
 $360$-unit chart.
 \item[Reduced Planck constant.] Action section whose regional coefficient
 and determinantal scale are constructed before the counterangle; in the
 ellipse, oriented area measure per unit of parametric phase.
 \item[Electron.] Persistent central section; its scalar readout composes
 additive--multiplicative incompatibility, regional transfer, closure defect
 and action return. The complete equation retains the section in addition
 to its mass value.
 \item[Barbero--Immirzi.] Oriented functional of angular transport and of
 returns in the hexad--octad incidence; in the constitutive collection, the
 coordinate that, together with normalized speed, reconstructs the labelled
 spectrum.
 \item[Catalan.] Endpoint value of the oriented depth-two moment of the
 dodecaphasic quarter-turn. Its collection is linked differentially to the
 vacuum response.
 \item[$G$.] Common coefficient of the radial realization relating
 energy and conjugate length while simultaneously preserving
 $H\Lambda=\hbar cI$ and $R\Lambda=L_\alpha^2I$.
 Longitudinal and polar composition determines its normalization;
 its dimensional expression is $c^3L_\alpha^2/\hbar$.
 \item[Boltzmann.] Transducer between the thermal line and decision energy,
 with declared tritic information normalization and thermal basis.
\end{description}
These intrinsic formulations rest on the constructed objects and equations.
Their scope is that of this realization; the physical identifications
retain their corresponding proofs and do not follow from structural
terminology alone.

\subsection{Transport of the nilpotent electronic block}
\label{md:compos:nilpotentes}

\paragraph{Proposition.}
Let
\begin{equation}
 S=\begin{pmatrix}1&0\\0&-1\end{pmatrix},\qquad
 J=\begin{pmatrix}0&1\\-1&0\end{pmatrix},\qquad
 n_\pm=\frac{S\pm J}{2},\qquad
 D=\operatorname{diag}(r,r^{-1}),\quad r>0.
 \label{md:compos:bloque-nilpotente}
\end{equation}
Define $n_{\pm,D}=Dn_\pm D^{-1}$ and $g=D^{-T}D^{-1}$. Then
\begin{equation}
 n_{\pm,D}^2=0,\qquad
 \{n_{+,D},n_{-,D}\}=I,\qquad
 (n_{+,D})^{\dagger_g}=n_{-,D},
 \label{md:compos:identidades-nilpotentes}
\end{equation}
where $A^{\dagger_g}=g^{-1}A^Tg$. The ellipse transported from
$Y^TY=2\hbar$ by $X=DY$ satisfies $X^TgX=2\hbar$.

\paragraph{Proof.}
Direct matrix calculation gives
$S^2=I$, $J^2=-I$ and $SJ+JS=0$. Consequently,
\[
 n_\pm^2=\tfrac14(S^2+J^2\pm SJ\pm JS)=0,
 \qquad
 n_+n_-+n_-n_+=\tfrac12(S^2-J^2)=I.
\]
Moreover, $S^T=S$ and $J^T=-J$ imply $n_+^T=n_-$. Conjugation by $D$
preserves products, squares and the identity. For the metric adjoint,
\[
 g^{-1}(Dn_+D^{-1})^Tg
 =DD^T D^{-T}n_-D^T D^{-T}D^{-1}
 =Dn_-D^{-1}=n_{-,D}.
\]
Finally, $D^TgD=I$, whence
$X^TgX=Y^TD^TgDY=Y^TY=2\hbar$. This proves the claims.

The exact algebraic verification takes place in
$\mathbb Q[r,r^{-1}]$; the restriction $r>0$ fixes the oriented elliptic
realization used here. This composition preserves the specific roles of the
different factors $1/2$ in the development: a numerical coincidence between
them does not identify their operators.
